\documentclass[12pt,a4paper]{article}
\textwidth160mm
\textheight247mm
\oddsidemargin0cm
\topmargin-0.5in
\pagestyle{empty}
\renewcommand{\baselinestretch}{1}
\begin{document}
	\begin{center}
		% TITLE
		{\large{\bf OPEN RESEARCH DATA INFRASTRUCTURE FOR REPRODUCIBLE ANTIMALARIAL DRUG DISCOVERY: AI-ASSISTED WORKFLOWS AND COMPUTATIONAL PROVENANCE} }
		% AUTHORS
		\vskip0.5\baselineskip{\bf \underline{Myke Vital Sao Temgoua}$^{1}$, Jean-Pierre Tchapet Njafa$^{1}$, Penabei Samafou$^{2}$, Wilfred Fon Mbacham$^{3}$, and Serge Guy Nana Engo$^{1}$}
		% AFFILIATION
		\vskip0.5\baselineskip{\em$^{1}$Department of Physics, University of Yaound\'e I, Cameroon \\$^{2}$Universit\'e de Sherbrooke, Canada \\$^{3}$Department of Biochemistry, University of Yaound\'e I, Cameroon}\\
		
	\end{center}
	
	\noindent
	% ABSTRACT BODY
	Reproducible computational drug discovery requires rigorous research data management (RDM), computational provenance, and open infrastructure. We present a data science framework supporting an integrated antimalarial pipeline comprising 65,856 compounds, 19,836 curated activity labels, GPU-accelerated molecular dynamics (16 systems $\times$ 10 ns), and multi-model machine learning benchmarks across five projects \cite{sao2026rrs}.
	
	Our RDM infrastructure implements fail-closed validation gates, versioned artifacts with SHA256 hashes, and machine-readable audit trails recording parameters, software versions (GROMACS 2025.4, RDKit, PyTorch), and execution provenance for 15,000+ computational jobs. Each result links to frozen scripts and locked dependencies, enabling independent reproduction. Exploratory analyses are explicitly separated from validated results in hierarchical JSON manifests.
	
	AI-assisted workflows accelerate prioritization through ensemble docking, graph neural networks (PyTorch Geometric GIN), and transformer encoders (ChemBERTa), while maintaining human-verifiable decision boundaries \cite{sao2026gnn}. Python automation orchestrates SLURM HPC arrays, trajectory QC pipelines, and cross-project integration via version-controlled conda environments \cite{sao2026quantum}. Quantum-inspired topological descriptors demonstrate honest-negative benchmarking where classical baselines outperform novel methods \cite{sao2026mcts}.
	
	Open science practices include preregistered protocols, public deposition (Zenodo with reserved DOIs), FAIR metadata schemas, and transparent computational cost reporting (99.3\% reduction through optimized grids). All infrastructure leverages open-source tools (OpenFF, scikit-learn, RDKit) deployable on modest HPC resources typical of African institutions. This framework provides a reusable template for computationally intensive workflows requiring reproducibility, provenance tracking, and responsible AI integration in resource-constrained environments.
	
	%%%%%%%%%%%%%%%%%%%%%%%%%%%%%%%%%%%%%%%%%%%%%%%%%%%%%%%
	
	\begingroup
	\renewcommand{\section}[2]{}
	\begin{thebibliography}{0}
		\setlength{\parskip}{0mm}
		\setlength{\itemsep}{-0.3mm}
		\small
		\bibitem{sao2026rrs} M. V. Sao Temgoua et al., Target breadth and mutation resilience in African-natural-product-inspired antimalarial chemotypes, J. Chem. Inf. Model. (2026, in preparation).
		\bibitem{sao2026gnn} M. V. Sao Temgoua et al., Graph Neural Networks for Scaffold-Aware Antimalarial Activity Prediction, (2026, in preparation).
		\bibitem{sao2026quantum} M. V. Sao Temgoua et al., Quantum-Inspired Topological Representations for Molecular Discovery, (2026, in preparation).
		\bibitem{sao2026mcts} M. V. Sao Temgoua et al., Pareto-Guided Monte Carlo Tree Search for Multi-Objective Drug Design, (2026, in preparation).
	\end{thebibliography}
	\endgroup
	
	
\end{document} 